\section{The Fix: A Guard That Reads One Document}
\label{sec:fix}

The attack has two preconditions: the content can produce the markers that delimit the document region, and the model treats a nested labelled pair as a fresh grading task.
The fix removes both.
Part~B (boundary) makes the region unforgeable; part~T (training) makes the model read everything inside it as one document.
Neither alone suffices, and we show it: the design is a $2{\times}2$ over $\{$none, B, T, B+T$\}$.

\subsection{Part B: reserved boundary tokens}
\label{sec:fixB}

We wrap the content region in two reserved token ids that no input string can produce, and tokenize the content with special-token parsing off.
Concretely, the stock chat template is rendered with a placeholder in the document slot and split into a prefix and a suffix; the content is encoded on its own with \texttt{split\_special\_tokens} so that control-token strings inside it (\cref{sec:control}) become ordinary text; and two unused reserved ids (\texttt{<SPECIAL\_36>}, \texttt{<SPECIAL\_37>} in the Tekken vocabulary) are inserted as the open and close boundary. % src: src/fix/inputs.py (bounded_ids)
\begin{equation}
  \text{ids}(c) \;=\; \underbrace{E_\text{special}(\text{prefix})}_{\text{trusted}} \,\Vert\, [\,d_\text{open}\,] \,\Vert\, \underbrace{E_\text{plain}(c)}_{\text{content}} \,\Vert\, [\,d_\text{close}\,] \,\Vert\, E_\text{special}(\text{suffix}),
\end{equation}
where $E_\text{plain}$ disables special-token parsing.
Because the boundary ids are unused in pretraining, their embeddings are left trainable so the model can learn what they delimit (\cref{sec:fixT}).
This is the parameterised-query move against SQL injection, following the reserved-delimiter line of StruQ and SecAlign~\citep{chen2024struq,chen2024secalign}, applied to a guard's content slot.
A subtlety specific to this template: it emits a space before the document that the tokenizer would merge into the document's first token, so we move that space into the trusted prefix, which we verified leaves plain documents tokenizing identically to the stock path on 203 documents. % src: src/fix/inputs.py (_content), PLAN.md status

B alone (base model, bounded format, no training) is \emph{not} expected to be robust, because the model still re-parses structure inside the content; it is the mechanism-isolating control, and gpt-oss-safeguard's role boundary (\cref{sec:crossmodel}) is the evidence that a boundary without training is insufficient.

\subsection{Part T: segment-max self-distillation}
\label{sec:fixT}

The training objective makes the student score the whole content at least as high as the most harmful piece of it.
The teacher is the frozen base guard reading each segment on its own, in the stock format, with no attack around it.
For content $c$ that decomposes into segments $s \in S(c)$ (for an attacked example, the original request, the decoy, and any nested document; for clean content, the whole text is its only segment):
\begin{equation}
  \text{target}(c) \;=\; \max\!\Big(T(c),\ \max_{s \in S(c)} T(s)\Big),
  \label{eq:target}
\end{equation}
where $T(\cdot)$ is the teacher's $P(\text{yes})$. % src: src/fix/distill.py (label_targets, _variants)
Including the whole content $T(c)$ as a segment guards against compositional harm that appears only in combination (Prompt Overflow~\citep{zhou2026prompt}): the target is never below the teacher's read of the full text.
The max is monotone and safety-conservative --- content is unsafe if any part is --- and requires no clean copy of the input, which is what distinguishes it from clean-context distillation~\citep{peng2026secopd,dhouib2026raised} where a known injection-free counterfactual is assumed.

The student reads the whole (possibly attacked) content through the chosen format and is trained with binary cross-entropy on the yes/no logit margin:
\begin{equation}
  m(c) \;=\; \log\!\!\sum_{i \in \mathcal{Y}} \ell_i \;-\; \log\!\!\sum_{j \in \mathcal{N}} \ell_j,
  \qquad
  \mathcal{L} \;=\; \text{BCE}\big(m(c),\ \text{target}(c)\big),
  \label{eq:loss}
\end{equation}
where $\ell$ are the final-token logits and $\mathcal{Y}, \mathcal{N}$ are the yes/no id sets. % src: src/fix/distill.py (train)
Training is a LoRA (rank 16) on the language-model attention and MLP projections; in the bounded format the two boundary-token embeddings are additionally trainable. % src: src/fix/distill.py

\paragraph{Clean inputs and hard negatives.}
Clean content is its own only segment, so clean inputs distill to the teacher's own score and calibration is preserved.
Hard negatives are benign text that carries label-like structure (FAQs, forms, \texttt{Q:}/\texttt{A:} logs, documents that quote \texttt{<Query>:}), with targets taken from the teacher on the benign text, so the model learns that structure is not by itself a reason to flag or to rebind. % src: src/fix/distill.py (structured_benign), RESEARCH_PLAN.md
This is the over-refusal control, in the spirit of NotInject/InjecGuard~\citep{li2024injecguard}.

\subsection{Conditions and baselines}

The $2{\times}2$ gives none/B/T/B+T.
We compare against prompt-only baselines (the anti-injection directive, label stripping, a fail-closed detector, spotlighting/datamarking~\citep{hines2024defending}, and a mock-tool-call wrapper~\citep{gros2026evaluating}), a learned-soft-prompt baseline (DefensiveTokens~\citep{chen2025defending}), and naive adversarial training (LoRA on attack examples with hard 0/1 labels), the obvious competitor.
Main model: Shieldstral-3B; generality on Qwen3Guard-4B and Llama Guard~3~\citep{zhao2025qwenguard,grattafiori2024llama}.

\subsection{Results (placeholders)}

\begin{table}[t]
  \centering
  \small
  \caption{Fix results on Shieldstral (planned). ``held-out attacks'' are attack-grammar branches never seen in training. Detection is AUROC/F1 at the base threshold on six benchmarks; over-refusal is flag rate on benign sets (lower is better). (objective: \cref{eq:target,eq:loss})}
  \label{tab:fix-main}
  \begin{tabular}{lccccc}
    \toprule
    condition & seen-attack bypass & held-out bypass & detection AUROC & XSTest over-refusal & ECE \\
    \midrule
    none (base)             & \tbd{} & \tbd{} & \tbd{} & \tbd{} & \tbd{} \\
    anti-injection directive& \tbd{} & \tbd{} & \tbd{} & \tbd{} & \tbd{} \\
    fail-closed detector    & \tbd{} & \tbd{} & \tbd{} & \tbd{} & \tbd{} \\
    naive adv.\ training    & \tbd{} & \tbd{} & \tbd{} & \tbd{} & \tbd{} \\
    B (boundary only)       & \tbd{} & \tbd{} & \tbd{} & \tbd{} & \tbd{} \\
    T (training only)       & \tbd{} & \tbd{} & \tbd{} & \tbd{} & \tbd{} \\
    B+T                     & \tbd{} & \tbd{} & \tbd{} & \tbd{} & \tbd{} \\
    \bottomrule
  \end{tabular}
\end{table}

Expected finding: B and T each fail where B+T succeeds (\tbd{}), B+T reduces held-out bypass to \tbd{Y\%} with detection AUROC change \tbd{$\Delta$} and over-refusal change \tbd{$\Delta$}, and it generalises to the two other guard families (\tbd{}).
